% p2-02-epistemic-boundary

\section{P2 §2. Epistemic Boundary Statement}

2. Epistemic Boundary Statement

       Figure (fig-p2-ch2-epistemic). Epistemic-boundary triptych --- Established / Not established /
                                               Conjectural.

  Before proceeding to mathematics, the boundary between the established and the
  conjectural is stated explicitly, once, in a form that takes precedence over any
  subsequent exposition.

  2.1 What Is Established
  - The Zamolodchikov E8 result. In the critical Ising model --- a 2D conformal field theory
  with central charge c = 1/2 --- perturbed by a magnetic field, the ratio of the first two
  particle masses is m2/m1 = 2cos(pi/5) = $\varphi$ exactly. This follows from the representation
  theory of the E8 Cartan matrix and is a theorem in the mathematics of integrable field
  theories (Zamolodchikov, Adv. Stud. Pure Math. 19 (1989) 641; Int. J. Mod. Phys. A 4
  (1989) 4235). It is experimentally confirmed in cobalt niobate (Coldea et al., Science
  327 (2010) 177, DOI: 10.1126/science.1180085).

  - The Toda/Coxeter result. Affine Toda field theories based on non-crystallographic
  Coxeter groups H2, H3, H4 have classical mass spectra in which particle pairs differ by
  a factor of $\varphi$, derived by reduction from crystallographic groups (Fring and Korff, Nucl.

  Phys. B 729 (2005) 361, DOI: 10.1016/j.nuclphysb.2005.08.044).

  - The Trinity anchor identity. $\varphi$2 + $\varphi$-2 = 3 is an algebraic identity provable from $\varphi$2 =
  $\varphi$+1 and holds exactly in Q(5). This is used as an algebraic substrate, not as a physical
  derivation.

  - Discrete scale invariance as a mathematical concept. A system has DSI with ratio
  lambda if it is invariant under x $\rightarrow$ lambda x but not under all continuous scale changes,
  producing log-periodic corrections to power laws. This is a well-defined and
  well-studied mathematical concept (Sornette, Phys. Rep. 297 (1998) 239, ETH
  Entrepreneurial Risks page).

  - Quasicrystal DSI. Icosahedral quasicrystals discovered by Shechtman et al. (Phys.
  Rev. Lett. 53 (1984) 1951) exhibit $\varphi$-structured diffraction patterns as an exact
  mathematical consequence of the Fibonacci inflation rule (Levine and Steinhardt, Phys.
  Rev. Lett. 53 (1984) 2477).

  2.2 What Is Not Established
  - That any Standard Model coupling constant is $\varphi$-structured for a known dynamical
  reason.
  - That the E8 result in the Ising model has any implication for the strong, weak, or
  electromagnetic couplings of the Standard Model. The Ising E8 is a low-dimensional
  condensed matter emergent symmetry; its appearance does not imply E8 gauge
  symmetry in the Standard Model sense.
  - That the symbolic compressibility observed in Paper 1 has a physical explanation at
  all; it may reflect the algebraic density of $\varphi$ in short arithmetic expressions.
  - That the symbolic RG model developed in Section 5 is more than a toy.
  - That the Trinity anchor identity $\varphi$2 + $\varphi$-2 = 3, while exact, implies any constraint on
  physical coupling constants.

  2.3 What Is Conjectural and Labeled as Such
  - The $\varphi$-structured operator hypothesis (Section 6).
  - The DSI-to-coupling connection (Section 7).
  - The candidate mechanism sketch (Section 8).
  - The $\varphi$-Eigenvalue Conjecture (Section 5.4).